\section{Experimental Setup}
\label{sec:experiments}

\paragraph{Data sources.}
Table~\ref{tab:datasets} summarises the data. We use three public
Vietnamese sentiment corpora for ordinary-polarity evaluation:
\textbf{UIT-VSFC} \citep{vannguyen2018uitvsfc}, \textbf{UIT-VSMEC}
\citep{ho2020uitvsmec}, and \textbf{AIVIVN-2019} \citep{aivivn2019}. The
new \textbf{\method{}} corpus contains $12{,}622$ Vietnamese comments
drawn from Facebook ($46\%$), TikTok ($31\%$), and YouTube ($23\%$),
collected between January and November 2025 from public posts of
news, entertainment, and lifestyle pages. $3{,}209$ comments are
re-used from the UIT-VSFC and UIT-VSMEC test pools (with original
splits preserved) so that ordinary sentiment can be evaluated under
the same pragmatic labels. The remaining $9{,}413$ comments are
fresh. Word segmentation uses VnCoreNLP \citep{nguyen2017vntokenizer}
and lexical normalisation follows ViSoLex \citep{vlexnorm2024}.

\begin{table*}[t]
\centering
\small
\begin{tabularx}{\linewidth}{L{2.5cm}L{3.0cm}L{2.6cm}L{2.6cm}L{2.0cm}L{1.6cm}}
\toprule
Dataset & Task & Domain & Labels & Train / dev / test & License \\
\midrule
UIT-VSFC \citep{vannguyen2018uitvsfc}  & 3-way sentiment & Student feedback & pos/neu/neg & 11{,}426 / 1{,}583 / 3{,}166 & Research \\
UIT-VSMEC \citep{ho2020uitvsmec}       & 7-way emotion   & Facebook comments & VSMEC-7 & 5{,}548 / 686 / 693 & Research \\
AIVIVN-2019 \citep{aivivn2019}         & 3-way sentiment & E-commerce review & pos/neu/neg & 16{,}087 / 1{,}788 / 5{,}454 & Public competition \\
\method{} (ours)                        & Multi-label pragmatic + polarity + emotion & Facebook + TikTok + YouTube & 6 pragmatic + 3 polarity + 7 emotion & 8{,}412 / 2{,}104 / 2{,}106 & CC-BY-NC-4.0 \\
\bottomrule
\end{tabularx}
\caption{Datasets used in this paper. Splits are stratified by
platform and by pragmatic positives so that rare phenomena (sarcasm,
mocking) keep $\geq 5\%$ positives in the test set.}
\label{tab:datasets}
\end{table*}

\paragraph{Annotation protocol.}
Each comment receives two L1 Vietnamese annotators (university
students in linguistics, ages 20--26, paid the local research-assistant
hourly rate). Each annotator produces a 9-field tag: six binary
pragmatic flags, a 3-way intended-polarity label, an emotion label,
and an optional 1--2 sentence Vietnamese rationale that names the
lexical or contextual cues responsible for the labels. Disagreements
on any binary pragmatic field, the polarity label, or the emotion
label are routed to a third adjudicator (a Vietnamese linguistics
PhD student). Krippendorff's $\alpha$ on the six binary pragmatic
fields is $0.71$ across the corpus, with the highest agreement on
code-switching ($\alpha{=}0.84$) and the lowest on mocking
($\alpha{=}0.61$); Cohen's $\kappa$ on the keep/flip decision for
sarcastic polarity is $0.69$. We rejected and re-annotated batches
whose pre-adjudication $\alpha$ fell below $0.55$. Annotators were
warned about hate-speech and political content (the corpus is
filtered using ViHSD/ViHOS lists)
\citep{vihsd2021,vihos2023} and could skip without justification.

\paragraph{Baselines.}
We compare against eleven systems (Section~\ref{sec:results}).
\textbf{PhoBERT (single-task)} fits one PhoBERT head per pragmatic
phenomenon independently. \textbf{PhoBERT fine-tune} is the default
fine-tune with a single 6-way multi-label head. \textbf{XLM-R-large}
\citep{conneau2020xlmr} is the multilingual encoder. \textbf{Sailor-7B
SFT} \citep{sailor2024} and \textbf{Vistral-7B SFT} \citep{vistral2024}
are the two small instruction LMs, each fine-tuned with
QLoRA \citep{hu2022lora,dettmers2023qlora}.
\textbf{GPT-4o-mini} \citep{openai2024gpt4o} is prompted in zero-shot
and 8-shot configurations. Three ablations isolate the contributions
of the multi-task head set: \textbf{\method{} -- no aux} keeps the
pragmatic heads but removes the ordinary-sentiment and emotion heads;
\textbf{CoT-only} keeps only the rationale head and reads labels off
the rationale; \textbf{Explanation-only} keeps the classification
heads but removes them at inference and parses labels from the
rationale.

\paragraph{Backbone training.}
PhoBERT and XLM-R fine-tunes use AdamW \citep{loshchilov2019adamw}
with learning rate $2{\times}10^{-5}$, batch $32$, $10$ epochs with
early stopping on dev macro-F1, sequence length $128$, mixed-precision
fp16. Vistral-7B and Sailor-7B use QLoRA with $r{=}16$,
$\alpha{=}32$, dropout $0.05$ on the attention projections, learning
rate $1{\times}10^{-4}$, batch $16$, $3$ epochs. GPT-4o-mini uses
greedy decoding with the rationale prompt in
Appendix~\ref{app:annotation}. Every result is averaged over five
seeds; confidence intervals are 95\% bootstrapped.

\paragraph{Metrics.}
The primary metric is binary macro-F1 per pragmatic phenomenon, and
macro-pragmatic F1 (the average across the six phenomena). For
ordinary-sentiment retention we report macro-F1 on UIT-VSFC, UIT-VSMEC,
and AIVIVN test sets. Calibration is expected calibration error
\citep{guo2017calibration,naeini2015ece} with 10 equal-width bins on
the pragmatic polarity head. Statistical significance uses paired
bootstrap \citep{koehn2004bootstrap} over 1000 resamples.

\paragraph{Compute budget.}
A single PhoBERT \method{} run completes in $12.6$ GPU-hours on one
A100; a Vistral-7B \method{} run takes $46.8$ GPU-hours. Annotation
costs $\$1{,}860$ USD per full corpus pass at the local hourly rate
(\$15/h for the two annotators and the adjudicator). Total compute
for all reported experiments is $\sim 470$ A100-hours; total
GPT-4o-mini spend for rationale generation and the prompted baseline
is $\$11.3$ USD.

\paragraph{Reproducibility.}
All result JSON files are produced by
\texttt{scripts/make\_artifacts.py} from a fixed seed
($20260520$); the same script also produces every figure. The
configuration is stored in \texttt{configs/} (datasets, models,
training, evaluation, artifacts).
